
Can a language model's architecture family be identified from its adversarial failures alone --- without knowing its weights, training data, or any other metadata? This paper investigates whether that answer is yes, at least in principle. Architecture family membership (encoder-only, decoder-only, encoder-decoder) accounts for approximately 29\% of variance in normalized adversarial vulnerability across attack categories at base model scale (permutation MANOVA $\eta^{2}=0.293$, Cohen's $f^{2}\approx 0.41)$, and this effect exceeds a pre-specified practical threshold in 83\% of the adversarial attack categories evaluated. The signal is large. However, statistical confirmation requires a larger model pool than the nine models studied here --- a limitation this paper quantifies as a concrete power analysis and study-design contribution.

The practical context motivates the question. A security auditor comparing encoder-only models (e.g., BERT-style) against decoder-only models (e.g., GPT-style) for adversarially sensitive deployment cannot currently quantify their different adversarial risk profiles from first principles. Robustness assessments today are model-specific and non-transferable: each new model requires a fresh evaluation from scratch. Architecture-family-level fingerprinting, if empirically confirmed, would allow principled inference about an unseen model's likely vulnerability profile from its architecture family membership alone, reducing evaluation cost substantially.

A deeper structural problem underlies evaluation cost. Existing adversarial robustness studies evaluate models as individuals, reporting scalar accuracy-drop metrics aggregated across attack types. Scalar aggregation discards the multivariate structure of adversarial vulnerability. Different attack types stress different model components: whether a model is more vulnerable to lexical substitutions than to paraphrase attacks may be systematically related to how its attention mechanism processes input perturbations. Encoder-only models with bidirectional attention globally redistribute attention mass over the full input sequence under perturbation; decoder-only models with causal attention process perturbations locally in sequence order. These structural differences, the present study hypothesizes, produce characteristic per-attack-category vulnerability profiles --- a $\Delta$*-vector --- that are consistent within architecture families and discriminable between them.

Prior work provides directional evidence. Neerudu et al. (2023) compare BERT, GPT-2, and T5 on GLUE perturbations and find GPT-2 (decoder-only) more robust than BERT (encoder-only) and T5 (encoder-decoder), a result consistent with architecture-family effects. Wang et al. (2021) establish AdvGLUE as a multi-task adversarial benchmark. Nie et al. (2020) provide ANLI as a human-crafted adversarial NLI benchmark. Zhang et al. (2026) find lower explanation flip rates for decoder LLMs in enterprise NLP contexts. Yet none of these works formalizes the $\Delta$*-vector representation, quantifies between-family effect size, or applies multivariate testing suited to small model pools. Sun et al. (2024, TrustLLM) evaluate 16 LLMs across trustworthiness dimensions but do not stratify by architecture family. Otmakhova et al. (2025, FLUKE) demonstrate that linguistic capability does not imply robustness to perturbations of that feature --- motivation consistent with the present approach but orthogonal in method.

This paper makes four contributions. First, it introduces the *\textit{$\Delta$}-vector framework\textbf{ for multivariate architecture-family adversarial profiling: each model is represented as a vector of normalized per-category vulnerability scores, and architecture-family effects are tested with permutation MANOVA. Second, it provides the }first formal effect-size measurement*\textit{ of between-family $\Delta$}-vector clustering: $\eta^{2}=0.293$ across nine models and six attack categories, exceeding the pre-specified $\eta^{2}=0.15$ threshold in five of six categories. Third, it contributes a \textbf{validated seven-module pipeline} (1,788 lines) covering fine-tuning via pre-existing checkpoints, adversarial evaluation, $\Delta$* computation, statistical analysis, and visualization. Fourth, it provides a \textbf{power analysis} establishing $N\geq 15$ ($\geq 5$ models per family) as the minimum for 80\% statistical power at the observed effect size --- concrete design guidance for a confirmatory follow-up study.

The paper proceeds as follows. Section 2 reviews adversarial robustness benchmarks, architecture comparison studies, and LLM trustworthiness evaluation. Section 3 describes the $\Delta$*-vector framework and statistical methods. Section 4 presents the experimental setup. Section 5 reports results. Section 6 discusses implications and limitations. Section 7 concludes.

